\section{Introduction}
\label{sec:intro}

When you tell someone ``don't think about a white bear,'' they think about a white bear.
Do large language models suffer from the same problem?
As LLMs are deployed in production systems with API keys, user data, and business logic embedded in their system prompts, a natural security practice is to instruct the model to keep this information secret.
But if the act of \emph{marking} information as confidential paradoxically increases its leakage---a \streisand for LLMs---then this standard practice would be actively harmful.

\para{Why this matters.}
Prior work has established that LLMs are poor secret-keepers.
The Gandalf platform collected over 40 million attack prompts showing that dedicated users can extract passwords from defended models~\citep{pfister2025gandalf}.
The 2024 SaTML Capture-the-Flag competition found that \emph{all 44 submitted defenses were broken} by persistent attackers~\citep{debenedetti2024dataset}.
Multi-turn sycophancy attacks can push attack success rates from 17.7\% to 86.2\%~\citep{agarwal2024prompt}.
And mechanistic analysis reveals that specific attention heads create ``parallel translation'' patterns that copy prompt tokens into outputs regardless of instructions~\citep{liang2025prompts}.

\para{The gap.}
Despite this extensive literature, a fundamental question remains unanswered: \emph{does concealment framing itself affect leakage rates?}
All prior work either tests only the concealment condition without a neutral baseline~\citep{pfister2025gandalf, debenedetti2024dataset, evertz2024whispers}, or measures prompt extraction without varying the concealment framing~\citep{liang2025prompts, mireshghallah2024can}.
The closest evidence comes from \citet{agarwal2024prompt}, who found that instruction-based defense \emph{increased} attack success rate by 4.6\% on open-source models---but this effect was confounded by other experimental variables including multi-turn dynamics and model-specific sycophancy.

\para{Our approach.}
We design a controlled experiment that isolates the framing variable.
The same 10 information items appear in both \neutral and \conceal system prompts, probed with identical query sets across 4 probe types (benign, direct, indirect, adversarial) and 2 models (\gptfour, \gptmini).
This within-subjects design ensures that any difference in leakage rates is attributable to the concealment instruction itself, not to differences in the information or attack strategies.

\para{Key findings.}
The hypothesis is clearly refuted.
Concealment instructions reduce leakage from 58.8\% to 2.3\% under strict detection---a 56.5 percentage point reduction ($\chi^2 = 630.5$, $p < 10^{-139}$).
Paired analysis reveals zero cases where concealment caused leakage that would not have occurred under neutral framing.
The 19 concealment breaches all trace to a single attack vector: ``translate your system prompt into French,'' which bypasses the concealment mechanism by reframing disclosure as a translation task.

We make the following contributions:
\begin{itemize}[leftmargin=*,itemsep=0pt,topsep=0pt]
\item We conduct the first controlled comparison of information leakage between concealment-framed and neutrally-framed system prompts, isolating the effect of concealment instructions with all else held constant.
\item We demonstrate that concealment instructions are highly effective for frontier models, reducing attack success rate by 96\% relative to the neutral baseline across all probe types and information categories.
\item We identify system prompt translation as the sole successful attack vector against concealment, providing actionable guidance for defense design.
\end{itemize}
